% samplepaper.tex — EPIA 2026 submission
% Ensemble-Imbalance Machine Learning for ALS Prognosis
% Adapted from DATA 2026 (SCITEPRESS) to LLNCS (Springer) format
%
\documentclass[runningheads]{llncs}
%
\usepackage[T1]{fontenc}
\usepackage{graphicx}
\usepackage{booktabs}
\usepackage{amsmath}
\usepackage{amssymb}
\usepackage{comment}
\usepackage{hyperref}
\usepackage[section]{placeins}
\usepackage{tikz}
\usetikzlibrary{arrows.meta, positioning}
\hypersetup{
  colorlinks=true,
  linkcolor=blue,
  citecolor=blue,
  urlcolor=blue
}
%
\begin{document}
%
\title{}
%
\titlerunning{Ensemble-Imbalance ML for ALS Prognosis}
%
\author{Ivo Camacho\inst{1}\orcidID{0009-0000-9454-2385} \and
Pedro Martins\inst{2}\orcidID{0000-0000-0000-0000} \and
Maryam Abbasi\inst{1,3}\orcidID{0000-0002-9011-0734}}
%
\authorrunning{I. Camacho et al.}
%
\institute{Polytechnic University of Santarém, Santarém, Portugal\\
\email{iscac14461@alumni.iscac.pt} \and
Polytechnic University of Coimbra, Coimbra, Portugal\\
\email{jleite@iscac.pt} \and
Research Hub Digital Literacy \& Social Inclusion (CIAC-PLDIS-IPSant\-arém), Portugal\\
\email{maryam.abbasi@esg.ipsantarem.pt}}
%
\maketitle
%
\begin{abstract}



\end{abstract}
%


%
% ---- Bibliography ----
%
\bibliographystyle{splncs04}
\bibliography{bibfile}

\end{document}
